\documentclass[11pt,a4paper]{article}
\usepackage[utf8]{inputenc}
\usepackage[T1]{fontenc}
\usepackage[french]{babel}
\usepackage{amsmath, amssymb}
\usepackage{geometry}
\geometry{margin=2.5cm}

\title{Correction du TD 1 : Analyse et Commande des systèmes linéaires continus}
\author{AU361}
\date{}

\begin{document}

\maketitle

\section*{Exercice 1 : Régulation de niveau d'un réservoir d'eau}

\textbf{1. Variables et paramètres du système} [cite: 516, 517, 518, 520, 521]\\
\textit{Paramètre :} La section du réservoir $S$ (en $m^2$).\\
\textit{Variables :}\\
- L'entrée (ou commande) $u(t)$ : Le débit entrant $q_e(t)$.\\
- La sortie à piloter $y(t)$ : Le niveau d'eau $n(t)$.\\
- La perturbation $\delta_1(t)$ : Le débit de fuite ou de consommation $q_s(t)$.

\textbf{2. Équation différentielle} [cite: 522]\\
La variation du volume d'eau dans le réservoir est égale à la différence entre le débit entrant et le débit sortant. Le volume étant $V(t) = S \cdot n(t)$, on a :
$$ \frac{dV(t)}{dt} = S \frac{dn(t)}{dt} = q_e(t) - q_s(t) $$

\textbf{3. Schéma fonctionnel du système asservi} [cite: 523, 524]\\
Le régulateur proportionnel $R$ produit la commande telle que $q_e(t) = K(cn(t) - n(t))$.\\
Le comparateur fait la différence entre $cn(t)$ (consigne) et $n(t)$ (mesure).\\
La perturbation $q_s(t)$ est soustraite au débit $q_e(t)$ avant d'attaquer le processus $H$.

\textbf{4. Équation différentielle en boucle fermée} [cite: 525]\\
En remplaçant $q_e(t)$ par son expression dans l'équation du système :
$$ S \frac{dn(t)}{dt} = K(cn(t) - n(t)) - q_s(t) \implies S \frac{dn(t)}{dt} + K n(t) = K cn(t) - q_s(t) $$

\textbf{5. Fonctions de transfert $R(p)$ et $H(p)$} [cite: 525, 526]\\
- Régulateur proportionnel : $R(p) = K$.\\
- Processus : En appliquant la transformée de Laplace avec conditions initiales nulles sur l'équation du réservoir $S p N(p) = Q_e(p) - Q_s(p)$. Si l'on ne considère que l'entrée $Q_e(p)$ (donc $Q_s=0$), on obtient $H(p) = \frac{N(p)}{Q_e(p)} = \frac{1}{S p}$.

\textbf{6. Diagramme de Bode de $H(p)$} [cite: 527]\\
La fonction $H(p) = \frac{1}{Sp}$ est un intégrateur pur.
- Gain : Droite de pente -20 dB/décade, passant par 0 dB à la pulsation $\omega = 1/S$.
- Phase : Constante et égale à $-90^\circ$.

\textbf{7. Fonctions de transfert en boucle fermée} [cite: 528, 529, 530]\\
A partir du schéma fonctionnel : $N(p) = H(p) [R(p)(CN(p) - N(p)) - \Delta_1(p)]$.
$$ N(p) = \frac{1}{Sp} [K(CN(p) - N(p)) - \Delta_1(p)] \implies N(p) \left(1 + \frac{K}{Sp}\right) = \frac{K}{Sp} CN(p) - \frac{1}{Sp} \Delta_1(p) $$
En multipliant par $\frac{Sp}{K}$, on identifie :
$$ T(p) = \frac{N(p)}{CN(p)} = \frac{1}{1 + \frac{S}{K}p} \quad \text{et} \quad T_{\delta 1}(p) = \frac{N(p)}{\Delta_1(p)} = \frac{-1/K}{1 + \frac{S}{K}p} $$

\textbf{8. Pôles et stabilité} [cite: 531]\\
Les deux fonctions de transfert ont le même dénominateur. Le pôle est solution de $1 + \frac{S}{K}p = 0$, soit $p = -\frac{K}{S}$.
Puisque $S > 0$ et $K > 0$, le pôle est strictement négatif. Le système asservi est donc **stable**.

\textbf{9. Diagrammes de Bode (comparaison)} [cite: 532]\\
- Boucle ouverte $bH(p) = R(p)H(p) = \frac{K}{Sp}$ : intégrateur pur (pente -20dB/décade, phase $-90^\circ$).
- Boucle fermée $T(p) = \frac{1}{1 + \frac{S}{K}p}$ : système du premier ordre (gain 0 dB en BF, puis cassure à $\omega_c = K/S$ avec pente -20 dB/décade ; phase de $0^\circ$ à $-90^\circ$).

\textbf{10. Régime permanent face à un échelon de perturbation} [cite: 533, 534]\\
Ici, $CN(p) = 0$ et $q_s(t)$ est un échelon unitaire $\implies Q_s(p) = \frac{1}{p}$.
$$ N(p) = T_{\delta 1}(p) Q_s(p) = \frac{-1/K}{1 + \frac{S}{K}p} \cdot \frac{1}{p} $$
En appliquant le théorème de la valeur finale :
$$ \lim_{t \to \infty} n(t) = \lim_{p \to 0} p N(p) = -\frac{1}{K} $$
Le niveau baisse de $1/K$ en régime permanent pour compenser la fuite.

\vspace{1cm}

\section*{Exercice 2 : Asservissement de vitesse d'une scie}

\textbf{1. Variables et paramètres} [cite: 538, 539, 540]\\
- Entrée : Tension de commande $u(t)$.
- Sortie : Vitesse angulaire $w(t)$.
- Perturbation : Couple résistant $Cr(t)$.
- Paramètres : $J, K_1, K_2, R$.

\textbf{2. Équation différentielle du système} [cite: 538, 541]\\
- Électrique : $u(t) = R i(t) + e(t) \implies i(t) = \frac{u(t) - K_2 w(t)}{R}$.
- Mécanique : $J \frac{dw(t)}{dt} = C_m(t) - Cr(t) = K_1 i(t) - Cr(t)$.
En combinant :
$$ J \frac{dw(t)}{dt} = \frac{K_1}{R}(u(t) - K_2 w(t)) - Cr(t) \implies J \frac{dw(t)}{dt} + \frac{K_1 K_2}{R} w(t) = \frac{K_1}{R} u(t) - Cr(t) $$

\textbf{3. Représentation d'état} [cite: 542]\\
On choisit la variable d'état $x(t) = w(t)$ et le vecteur d'entrée $V(t) = \begin{bmatrix} u(t) \\ Cr(t) \end{bmatrix}$.
$$ \dot{x}(t) = -\frac{K_1 K_2}{R J} x(t) + \begin{bmatrix} \frac{K_1}{R J} & -\frac{1}{J} \end{bmatrix} \begin{bmatrix} u(t) \\ Cr(t) \end{bmatrix} $$
$$ y(t) = x(t) $$

\textbf{4. Fonction de transfert $T(p)$} [cite: 543]\\
Si $Cr(t) = 0$, la transformée de Laplace donne $(J p + \frac{K_1 K_2}{R}) W(p) = \frac{K_1}{R} U(p)$.
$$ T(p) = \frac{W(p)}{U(p)} = \frac{K_1/R}{Jp + \frac{K_1 K_2}{R}} = \frac{1/K_2}{1 + \frac{RJ}{K_1 K_2}p} $$

\textbf{5. Stabilité} [cite: 544]\\
Le pôle est $p = -\frac{K_1 K_2}{R J}$. Comme toutes les grandeurs physiques sont positives, le pôle est réel strictement négatif. Le système est stable.

\textbf{6. Application d'un échelon $u(t)=10V$} [cite: 545, 546, 547]\\
- En régime permanent ($t \to \infty$) : $\lim_{t \to \infty} w(t) = \lim_{p \to 0} p \cdot T(p) \frac{10}{p} = \frac{10}{K_2}$.
- A $t=0^+$ : La variable d'état $w(t)$ est continue. Etant initialement nulle, $w(0^+) = 0$.
- Accélération à $t=0^+$ : À partir de l'équation différentielle, à $t=0^+$, $w(0^+)=0$ donc $J \cdot acc(0^+) + 0 = \frac{K_1}{R} \cdot 10 \implies acc(0^+) = \frac{10 K_1}{R J}$.

\textbf{7. Prise en compte de la perturbation} [cite: 548, 549, 550]\\
L'équation différentielle complète dans le domaine de Laplace est :
$$ (J p + \frac{K_1 K_2}{R}) W(p) = \frac{K_1}{R} U(p) - CR(p) $$
$$ W(p) = \frac{1/K_2}{1 + \frac{R J}{K_1 K_2} p} U(p) - \frac{R/(K_1 K_2)}{1 + \frac{R J}{K_1 K_2} p} CR(p) $$
On identifie $T_\delta(p) = \frac{-R/(K_1 K_2)}{1 + \frac{R J}{K_1 K_2} p}$.

\vspace{1cm}

\section*{Exercice 3 : Régulation de température d'une enceinte}

\textbf{1. Modélisation} [cite: 565, 567, 568]\\
- Paramètres : $\alpha_1, \alpha_2, \alpha_3$.
- Variables : $Tint(t), Text(t), Tpar(t), Puis(t)$.
- Commande : $Puis(t)$ ; Sortie : $Tint(t)$ ; Perturbation : $Text(t)$.

\textbf{2. Représentation externe} [cite: 570, 571, 572, 573]\\
Dans le domaine de Laplace :
(1) $p \cdot Tpar = \alpha_1(Tint - Tpar) + \alpha_2(Text - Tpar) + Puis \implies (p + \alpha_1 + \alpha_2)Tpar = \alpha_1 Tint + \alpha_2 Text + Puis$
(2) $p \cdot Tint = \alpha_3(Tpar - Tint) \implies Tpar = \frac{p+\alpha_3}{\alpha_3} Tint$.
En injectant (2) dans (1), on trouve :
$$ [p^2 + p(\alpha_1 + \alpha_2 + \alpha_3) + \alpha_2 \alpha_3] Tint = \alpha_3 Puis + \alpha_2 \alpha_3 Text $$
Donc :
$$ H(p) = \frac{\alpha_3}{p^2 + p(\alpha_1+\alpha_2+\alpha_3) + \alpha_2 \alpha_3} \quad \text{et} \quad Hd(p) = \frac{\alpha_2 \alpha_3}{p^2 + p(\alpha_1+\alpha_2+\alpha_3) + \alpha_2 \alpha_3} $$
Gain statique de $Hd(p)$ (pour $p=0$) : $K_{Hd} = \frac{\alpha_2 \alpha_3}{\alpha_2 \alpha_3} = 1$.

\textbf{3. Enceinte parfaitement isolée ($\alpha_2 = 0$)} [cite: 575, 576, 577, 578, 579]\\
- a) La fonction $Hd(p)$ devient nulle, l'extérieur n'a plus d'influence.
- b) $H(p) = \frac{\alpha_3}{p^2 + p(\alpha_1+\alpha_3)} = \frac{\alpha_3}{p(p + \alpha_1 + \alpha_3)} = \frac{\frac{\alpha_3}{\alpha_1+\alpha_3}}{p(1 + \frac{1}{\alpha_1+\alpha_3} p)}$. On a bien la forme demandée avec $G = \frac{\alpha_3}{\alpha_1+\alpha_3}$ et $T = \frac{1}{\alpha_1+\alpha_3}$.
- c) La réponse indicielle s'obtient en intégrant un système du 1er ordre. C'est une rampe amortie : $Tint(t) = G \cdot [t - T(1 - e^{-t/T})]$.

\textbf{4. Enceinte avec pertes thermiques} [cite: 581, 583, 584, 585]\\
- On calcule le dénominateur avec $\alpha_1=0.9, \alpha_2=0.01, \alpha_3=0.1$ : $p^2 + 1.01 p + 0.001$.
- Ce polynôme admet deux racines : $p_1 \approx -0.001$ et $p_2 \approx -1$. Le pôle $p_1$ est dominant (constante de temps très grande), ce qui explique le comportement apparent du premier ordre.
- Sur la Figure 2, on atteint $63\%$ de la valeur finale vers $t = 1000$s, d'où $\tau = 1000$s.
- Le gain statique de $H(p)$ est $K_H = \frac{0.1}{0.001} = 100$. Le gain statique de $Hd(p)$ est 1.

\vspace{1cm}

\section*{Exercice 4 : Régulation de température d'un cryostat}

\textbf{Identification des paramètres du modèle} [cite: 614, 619, 620, 623, 627, 630, 634, 635, 645, 650]\\
L'essai indique que $Text(t) = 300K$ constant.
Le système est à l'équilibre avant $t=100s$ (où $U=0$). $Temp(t<100) = Hd(0) \cdot 300 = 300 K$ (puisque $den(0)=1$). Cela correspond au graphe.

À $t=100s$, un échelon de consigne de vitesse est appliqué : passage à 1000 tr/min.
Sachant que 30V $\leftrightarrow$ rapport cyclique unitaire $U=1 \leftrightarrow$ 3000 tr/min, la vitesse est $3000 \cdot U(t)$.
Pour 1000 tr/min, l'échelon de commande est donc $\Delta U = \frac{1000}{3000} = \frac{1}{3}$.

Sur le graphe, la température passe de 300 K à 200 K en régime permanent. L'amplitude de la variation de la sortie est $\Delta Temp = -100 K$.
L'équation en régime permanent pour cette variation donne : $\Delta Temp = H(0) \cdot \Delta U = \frac{G}{den(0)} \cdot \Delta U$.
Ainsi, $-100 = G \cdot \frac{1}{3} \implies G = -300$.

Pour le dénominateur, la réponse a l'allure d'un premier ordre.
La valeur finale de la variation est 200K. À $63\%$ de la transition (soit une chute de 63K, donc Temp=237K), on lit sur le graphe un temps d'environ $300s$. Sachant que l'échelon a démarré à $100s$, la constante de temps est $\tau = 300 - 100 = 200s$.

Finalement, $den(p)$ pour un système du 1er ordre s'écrit $1 + \tau p$.
Les paramètres sont donc :
$$ G = -300 \quad \text{et} \quad den(p) = 1 + 200p $$

\end{document}